\documentclass[11pt,a4paper]{article}
\usepackage[margin=1in]{geometry}
\usepackage{enumitem}
\usepackage[hidelinks]{hyperref}
\usepackage{titlesec}
\usepackage{setspace}

% Formato general
\setlength{\parskip}{0.5em}   % espacio entre párrafos
\setlength{\parindent}{1.25em} % indentación primer renglón

% Para títulos, sin indentación y con espacio inferior
\titleformat{\section}{\large\bfseries}{}{0em}{}
\titlespacing*{\section}{0pt}{1em}{0.5em} % margen izq, espacio arriba y abajo

% Unificar sangría y espacio para listas itemize
\setlist[itemize]{leftmargin=2em, noitemsep, topsep=0.18em}

% Definir un entorno general para sangría uniforme
\newenvironment{uniformindent}{%
  \begin{list}{}{%
    \setlength{\leftmargin}{2em}
    \setlength{\rightmargin}{0pt}
    \setlength{\topsep}{0pt}
    \setlength{\partopsep}{0pt}
    \setlength{\parsep}{0pt}
    \setlength{\itemsep}{0pt}
  }
  \item[]
}{%
  \end{list}
}
\newenvironment{indentedpar}{%
  \begin{list}{}{%
    \setlength{\leftmargin}{2em}
    \setlength{\rightmargin}{0pt}
    \setlength{\itemsep}{0pt}
    \setlength{\parsep}{0pt}
    \setlength{\topsep}{0pt}
  }
  \item[]
}{%
  \end{list}
}
\begin{document}

\begin{center}
    {\LARGE \textbf{Alejandro Chavez}} \\[0.15cm]
    Ingeniero en Ciencias de la Computaci\'on \\[0.2cm]
    \begin{tabular}{@{} c @{}}
        \href{mailto:chavezalejo85@gmail.com}{chavezalejo85@gmail.com} \quad | \quad +593 999274384 \quad | \quad \href{https://www.linkedin.com/in/alech7/}{linkedin.com/in/alech7}
    \end{tabular}
\end{center}

\begin{uniformindent}

\section*{Educaci\'on}

\begin{samepage}
\textbf{Maestr\'ia en Inteligencia Artificial} \hfill Agosto 2026 -- Agosto 2027\\
Universidad San Francisco de Quito\\
\end{samepage}

\begin{samepage}
\textbf{Ingenier\'ia en Ciencias de la Computaci\'on} \hfill Septiembre 2021 -- Febrero 2026\\
Escuela Polit\'ecnica Nacional\\
\end{samepage}

\section*{Experiencia}

\textbf{Freelancer} \hfill Abril 2026 -- Presente
\begin{itemize}
    \item Desarrollo de servicios backend con .NET (C\#) y aplicaciones m\'oviles en Swift para clientes.
    \item Creaci\'on de un bot para agenda de citas en lenguaje natural usando n8n, WhatsApp Business y Next.js.
    \item Creaci\'on de un sistema RAG para la recuperaci\'on de memorias del agente de IA por usuario.
    \item Integraci\'on de contacto directo con el doctor dentro del flujo de agendamiento.
\end{itemize}

\textbf{Busmatick Group} \hfill Abril 2026 -- Agosto 2026
\begin{itemize}
    \item Soporte a usuarios y al sistema.
    \item Revisi\'on de logs, bases de datos y cuentas de usuarios para resolver incidentes y apoyar la operaci\'on.
    \item Automatizaci\'on de procesos para mejorar flujos de soporte y confiabilidad del sistema.
\end{itemize}

\textbf{Ecuavisa} \hfill Agosto 2025 -- Febrero 2026
\begin{itemize}
    \item Levantamiento de procesos y automatizaci\'on con IA (LLMs, visi\'on por computador y ML).
    \item Construcci\'on de soluciones con Google Cloud, la plataforma de OpenAI, Apache Spark y Kafka.
\end{itemize}

\textbf{BlueHat Corp} \hfill Marzo 2025 -- Agosto 2025
\begin{itemize}
    \item Ejecuci\'on de pruebas de ingenier\'ia social, OSINT y pentesting web para evaluar exposici\'on, riesgos reputacionales y postura de seguridad.
    \item Apoyo en iniciativas de gobierno de ciberseguridad alineadas con CGCI, incluyendo seguimiento con base en ISO/IEC 27001.
\end{itemize}

\textbf{Pasant\'ia en SLB} \hfill Octubre 2024 -- Marzo 2025
\begin{itemize}
    \item Pasante desarrollador en Power Platform, con dise\~no y desarrollo de soluciones en PowerApps para optimizar el seguimiento del consumo de di\'esel y la gesti\'on de activos.
    \item Implementaci\'on de formularios y reglas de negocio para validaci\'on de datos y mejora de la calidad y exactitud de la informaci\'on registrada.
    \item Contribuci\'on a modelos de datos en Dataverse, componentes UX/UI y gesti\'on \'agil del proyecto mediante Azure DevOps.
\end{itemize}

\textbf{Pasant\'ia en IQVIA} \hfill Marzo 2024 -- Mayo 2024
\begin{itemize}
    \item Pasant\'ia en soporte y administraci\'on de bases de datos, con automatizaci\'on y verificaci\'on de procesos semanales mediante ETL, SQL y PowerShell.
    \item Dise\~no y desarrollo de paneles en Tableau y Power BI para visualizaci\'on y seguimiento de KPIs.
    \item Dise\~no y desarrollo de consultas personalizadas para atender necesidades espec\'ificas de clientes y garantizar la entrega precisa de datos.
\end{itemize}

\section*{Proyectos Personales}

\textbf{Busqua} \hfill \href{https://busqua.com}{busqua.com}
\begin{indentedpar}
Plataforma de b\'usqueda de productos en tiempo real que conecta proveedores, ubicaci\'on y datos de contacto. Tecnolog\'ias: AWS, Odoo, Elixir, Python, SQL y React.
\end{indentedpar}

\textbf{Retrieval Augmented System}
\begin{indentedpar}
Desarroll\'e un sistema RAG para recuperar datos de un corpus personalizado y generar res\'umenes concisos y contextualizados.
\end{indentedpar}

\section*{Certificaciones}

\textbf{AWS Academy Cloud Foundations}\\
Amazon Web Services (AWS) \\ Emitido: Julio 2025\\

\textbf{Network Defense Essentials}\\
EC-Council \\ Emitido: Enero 2025 -- Expira: Enero 2028\\
ID de credencial: ECC5837692041

\section*{Publicaciones}

\textit{``Dise\~no y evaluaci\'on de un pipeline de b\'usqueda sem\'antica basado en RAG para la recuperaci\'on de informaci\'on cient\'ifica''} --- Publicaci\'on pendiente.\\
Revisi\'on sistem\'atica sobre los fundamentos de los sistemas RAG, sus componentes clave, flujos de implementaci\'on, aplicaciones en investigaci\'on cient\'ifica y desaf\'ios actuales.

\section*{Idiomas y Habilidades}

\textbf{Programaci\'on:} Python, Java, JavaScript, C\# (.NET backend), Swift, HTML, CSS\\
\textbf{Herramientas:} Kali Linux, SQL, React, Azure, PowerShell\\
\textbf{Metodolog\'ias:} Scrum, XP\\
\textbf{Otros:} Ingl\'es fluido, franc\'es b\'asico, conocimientos en sistemas Cisco

\end{uniformindent}

\end{document}
